% ============================================================================
% SEC03.TEX - Section 12.3: Menggunakan Pool
% ============================================================================

\section{Menggunakan Pool}

Sekarang ubah cara Player menembak.

\subsection{Update PlayerController}
Ganti \texttt{Instantiate} dengan mengambil dari Pool.

\begin{lstlisting}[language={[Sharp]C}]
void Shoot()
{
    // Minta peluru dari Pool
    GameObject bullet = BulletPool.Instance.GetBullet();
    
    if (bullet != null)
    {
        bullet.transform.position = firePoint.position;
        bullet.transform.rotation = firePoint.rotation;
        
        // Reset kecepatan peluru (penting!)
        // ... (Logika reset ada di BulletController)
    }
}
\end{lstlisting}

\subsection{Update BulletController}
Jangan \texttt{Destroy}! Tapi non-aktifkan.

\begin{lstlisting}[language={[Sharp]C}]
void OnEnable()
{
    // Timer manual untuk mematikan peluru setelah 2 detik
    Invoke("DisableBullet", 2f);
}

void DisableBullet()
{
    gameObject.SetActive(false); // Kembalikan ke pool
}

private void OnTriggerEnter2D(Collider2D other)
{
    if (other.CompareTag("Enemy"))
    {
        DisableBullet(); // Disable, bukan Destroy
        // Destroy(other.gameObject); // Asteroid tetap destroy (atau dipool juga)
    }
}

void OnDisable()
{
    CancelInvoke(); // Bersihkan timer
}
\end{lstlisting}
